\documentclass[12pt]{exam}
\usepackage[utf8]{inputenc} 
\usepackage[brazil]{babel}
\usepackage{graphicx}
\usepackage{latexsym}
\usepackage{listings}
\usepackage{amsmath}
\usepackage{multirow}
\usepackage{multicol}
\usepackage{amsfonts}
\usepackage{amssymb}
\usepackage{epsfig,psfrag} 
\usepackage{wrapfig}
\usepackage[left=1.8cm,top=1cm,right=1.8cm,bottom=1cm]{geometry}
\usepackage[lined, ruled, vlined]{algorithm2e}
\usepackage[normalem]{ulem}
\usepackage{graphics}
\usepackage{pstricks,pst-node,pst-tree}
\usepackage{eqnarray,amsmath}
\usepackage{caption}
\usepackage{circuitikz}
\usepackage{datetime}
\usepackage{nicematrix, tikz}
\usepackage{mathtools}
\ctikzset{bipoles/thickness=1}
\usepackage[colorlinks=true, allcolors=blue]{hyperref}
\usepackage{listings}
\usepackage{xcolor}
\usepackage{ragged2e}

\definecolor{cinza-codigo}{gray}{0.92}
\setlength{\fboxsep}{1.2pt}

\ExplSyntaxOn
\cs_new_protected:Npn \__code_word:n #1 {
  \tl_if_empty:nF { #1 }
   {
     \colorbox{cinza-codigo}{\texttt{#1}}
   }
}

\NewDocumentCommand{\code}{m}
 {
   \seq_set_split:Nnn \l_tmpa_seq { ~ } { #1 }
   \int_set:Nn \l_tmpa_int { 0 }
   \seq_map_inline:Nn \l_tmpa_seq
    {
      \int_incr:N \l_tmpa_int
      \int_compare:nF { \l_tmpa_int == 1 } { \allowbreak\ \hspace{0pt} }
      \tl_set:Nn \l_tmpa_tl { ##1 }
      \tl_replace_all:Nnn \l_tmpa_tl { \_ } { \__code_word:n { \_ } \allowbreak }
      \tl_replace_all:Nnn \l_tmpa_tl { ( } { \__code_word:n { ( } \allowbreak }
      \tl_replace_all:Nnn \l_tmpa_tl { ) } { \__code_word:n { ) } \allowbreak }
      \tl_replace_all:Nnn \l_tmpa_tl { [ } { \__code_word:n { [ } \allowbreak }
      \tl_replace_all:Nnn \l_tmpa_tl { ] } { \__code_word:n { ] } \allowbreak }
      \tl_replace_all:Nnn \l_tmpa_tl { , } { \__code_word:n { , } \allowbreak }
      \tl_replace_all:Nnn \l_tmpa_tl { : } { \__code_word:n { : } \allowbreak }
      \tl_replace_all:Nnn \l_tmpa_tl { / } { \__code_word:n { / } \allowbreak }
      \tl_replace_all:Nnn \l_tmpa_tl { . } { \__code_word:n { . } \allowbreak }
      \__code_word:n { \l_tmpa_tl }
    }
 }
\ExplSyntaxOff

% 2. Configurar o estilo do bloco de código
\lstset{
    language=bash,                  % Define a linguagem como Bash
    basicstyle=\ttfamily\small,     % Fonte monoespaçada e menor
    keywordstyle=\color{blue},      % Cor para comandos/palavras-chave
    commentstyle=\color{gray},      % Cor para comentários
    stringstyle=\color{purple},     % Cor para strings
    numbers=left,                   % Mostra números de linha na esquerda
    numberstyle=\tiny\color{gray},  % Estilo dos números de linha
    stepnumber=1,                   % Número em cada linha
    backgroundcolor=\color{lightgray!20}, % Cor de fundo do bloco
    frame=single,                   % Moldura ao redor do código
    breaklines=true,                % Quebra linhas longas automaticamente
    showstringspaces=false          % Não mostra espaços especiais em strings
}

\newcommand\encircle[1]{%
  \tikz[baseline=(X.base)] 
    \node (X) [draw, shape=circle, inner sep=0] {\strut #1};}

\newdateformat{specialdate}{\THEYEAR-\twodigit{\THEMONTH}-\twodigit{\THEDAY}}
\newcommand{\examdate}{Data de Entrega: 11/12/2024}

%\pagestyle{headandfoot}
%\firstpagefooter{}{Pág. \thepage\ de \numpages}{}
%\firstpageheader{}{}{}
%\runningfooter{}{Pág. \thepage\ de \numpages}{}
%\runningheader{}{}{}

\extraheadheight[.25in]{.5in}
\extrafootheight[.25in]{.5in}
\pagestyle{empty}
\usepackage{xcolor}
\definecolor{greendel}{HTML}{2F5416}
\definecolor{bluepoli}{HTML}{1E1D44}
\definecolor{redpoli}{HTML}{C73834}




\begin{document}
%\setlength{\arrayrulewidth}{1cm}
\setlength{\arrayrulewidth}{1.5pt}
%\centerline{\rule{\textwidth}{0.2mm} \\}\\
\vspace{0.1cm}
%\begin{tabular}{l!{\color{bluepoli}\vline}l}
\begin{tabular}{l|l}

\thispagestyle{empty}
\multirow{4}{30mm}{\vspace*{-1cm}\hspace*{-0.4cm}\includegraphics[width=0.15\textwidth]{logos/ciaa_gray.pdf}}
& {\bf \Large{Revisão Acelerada de Programação - RAP}} \\
 & {\bf \Large{RAP-2026, Marinha do Brasil}} \\
 & {Centro de Instrução Almirante Alexandrino}\\
 & {Departamento de Aperfeiçoamento Avançado para Oficiais}\\
 & {Curso de Aperfeiçoamento Avançado em Guerra Acústica (C-ApA-Of-GA)}\\
 & {Professor: João Victor da Fonseca Pinto}\\
 & {Descrição do Documento: Planejamento de aulas}\\
 & {}\\

%\hspace{0.0cm} & {Data: \today}\\
\end{tabular}\\ 
%\centerline{\rule{\textwidth}{0.2mm}}


\thispagestyle{headandfoot}
\normalsize
\vspace{-0.2cm}

\vspace{0.2cm}




\section*{Introdução}


Esta disciplina de Revisão Acelerada de Programação em Python tem por finalidade consolidar os conceitos fundamentais de lógica de programação e desenvolvimento de algoritmos sob a perspectiva da Guerra Acústica. O conteúdo é estruturado de forma intensiva, visando capacitar os militares na manipulação de dados complexos, na automação de rotinas de análise e na implementação rápida de soluções computacionais aplicadas. Mediante uma abordagem prática e direcionada, os alunos serão instrumentalizados a utilizar a linguagem Python como ferramenta de suporte técnico-operacional, otimizando o processamento de fluxos de informação e garantindo maior celeridade no cumprimento das tarefas correlatas à especialidade.


\section*{Planejamento de Aulas}

Conforme o cronograma estabelecido pela organização do curso, as aulas presenciais estão destacadas em vermelho, e as telepresenciais, em azul.

\begin{itemize}\justifying
    \item \textbf{Aula 1 (\textcolor{red}{27 de julho de 2026}) - 3 horas:}
    \begin{itemize}\justifying
        \item Apresentação da matéria: Método de avaliação e cronograma.
        \item Conceito de Python: Entender o que é Python e por que ele é amplamente utilizado.
        \item Conceito de Docker: Compreender o que são contêineres e sua importância para o desenvolvimento robusto.
        \item Conceito de Docker Desktop: Conhecer a interface que gerencia nossos contêineres.
        \item Conceito de Jupyter Lab: Aprender a utilizar e navegar na interface de notebooks.
        \item Conceito de IDE: Aprender a utilizar o Microsoft VS Code.
        \item Primeiros Passos: Aprender atalhos importantes e executar seu primeiro código em Python: no terminal, em um script e em um Jupyter notebook.        
    \end{itemize}
    
    \item \textbf{Aula 2 (\textcolor{blue}{3 de agosto de 2026} e \textcolor{blue}{10 de agosto de 2026)} - 5 horas:}
    \begin{itemize}\justifying
        \item Tipos e Variáveis: Entender como criar variáveis e como o Python categoriza os dados básicos.
        \item Casting: Aprender a converter dados de um tipo para outro.
        \item Built-in Functions: Explorar funções integradas como \code{print}, \code{help}, \code{len}, \code{abs}, etc.
        \item Tipos de Prints e Formatações: Dominar técnicas para exibir resultados formatados (como f-strings).
        \item Estruturas Condicionais: Tomar decisões no código usando \code{if}, \code{elif} e \code{else}.
        \item Introdução a Funções: Entender o que é uma função, para que serve e como criar as suas próprias.
        \item Exemplos Introdutórios: Praticar com mini-programas que unem todos os conceitos.
    \end{itemize}
    
    \item \textbf{Aula 3 (\textcolor{red}{17 de agosto de 2026}) - 3 horas:}
    \begin{itemize}\justifying
        \item Estruturas de Repetição: Entender os conceitos de loops (\code{for} e \code{while}).
        \item Repetição com \code{for}: Dominar a iteração sobre sequências e a utilização da função \code{range()} dentro de funções.
        \item Repetição com \code{while}: Compreender como repetir instruções baseadas em condições lógicas variáveis.
        \item Controle de Fluxo com \code{break} e \code{continue}: Aprender a interromper ou pular iterações de loops de forma inteligente.
        \item Exercícios Práticos: Aplicar essas técnicas em exercícios práticos.
    \end{itemize}


    \item \textbf{Aula 4 (\textcolor{blue}{24 de agosto de 2026} e \textcolor{blue}{28 de agosto de 2026}) - 5 horas:}

    \begin{itemize}\justifying
        \item Listas em Python: Compreender o conceito de listas, declaração de listas vazias e preenchidas, e indexação positiva/negativa.
        
        \item Principais Métodos e Funções: Dominar funções como \code{append()}, \code{insert()}, \code{remove()}, \code{pop()}, \code{len()}, \code{sort()} e \code{reverse()}.

        \item Listas de Listas (Matrizes): Entender o conceito de listas aninhadas, acesso bidimensional a elementos e iteração em matrizes.

        \item Exercícios Práticos: Aplicar essas técnicas em exercícios práticos.

    \end{itemize}


    \item \textbf{Aula 5 (\textcolor{blue}{9 de setembro de 2026}) - 3 horas:}
    \begin{itemize}\justifying
        \item Conceito de String: Entender a representação de texto, indexação, fatiamento e imutabilidade.
        \item Operações Básicas: Dominar concatenação, repetição, tamanho (\code{len()}) e operadores de pertinência (\code{in}).
        \item Métodos de Modificação: Aprender \code{.upper()}, \code{.lower()}, \code{.strip()} e \code{.replace()}.
        \item Divisão e Junção: Compreender \code{.split()} e \code{.join()} para converter entre texto e listas.
        \item Busca e Validação: Utilizar \code{.find()}, \code{.startswith()} e \code{.endswith()}.
        \item Exercícios Práticos: Aplicar essas técnicas em exercícios práticos.
    \end{itemize}



    \item \textbf{Aula 6 (\textcolor{blue}{11 de setembro de 2026}) - 3 horas:} 
    \begin{itemize}\justifying
       \item Conceito de Dicionário: Compreender o que é um dicionário, como chaves e valores são associados e suas propriedades básicas.
       \item Métodos e Funções: Dominar o uso de \code{.get()}, \code{.keys()}, \code{.values()}, \code{.items()}, \code{.update()}, \code{.pop()} e a palavra-chave \code{del}.
       \item Vantagem em Buscas: Demonstrar na prática a diferença de eficiência (complexidade de busca) entre listas e dicionários.
      \item Exercícios Práticos: Aplicar essas técnicas em exercícios práticos.

   \end{itemize}

    \item \textbf{Aula 7 (\textcolor{blue}{14 de setembro de 2026}) - 3 horas:}
    \begin{itemize}\justifying
        \item Mutabilidade vs Imutabilidade: Compreender a diferença conceitual entre objetos que podem ser alterados na memória e aqueles que são protegidos contra escrita, especialmente ao serem passados para funções.
        \item Conceito de Tupla: Compreender o que é uma tupla, como ela se diferencia de listas e dicionários.
        \item Casos de Uso: Entender para que servem as tuplas (garantia de integridade, desempacotamento de dados e chaves de dicionários).
        \item Métodos de Tuplas: Conhecer a tabela de métodos específicos de tuplas (\code{.count()} e \code{.index()}).
        \item Exemplos e Unpacking: Dominar a sintaxe de declaração, fatiamento e desempacotamento.
        \item Exercícios Práticos: Aplicar essas técnicas em exercícios práticos.
        \item Disponibilização da primeira lista de exercícios.
    \end{itemize}

    \item \textbf{Aula 8 (\textcolor{blue}{18 de setembro de 2026}) - 2 horas:}
    \begin{itemize}\justifying
        \item Entender a Abertura de Arquivos: Aprender sobre a função \code{open()} e seus diferentes modos (\code{'r'}, \code{'w'}, \code{'a'}).
        \item Utilizar o Gerenciador de Contexto: Aprender a melhor prática de abertura com a estrutura \code{with}.
        \item Escrever em Arquivos: Criar novos arquivos de    texto e adicionar linhas com \code{write()} e \code{writelines()}.
        \item Ler Arquivos: Utilizar métodos de leitura como \code{read()}, \code{readline()}, \code{readlines()} e loops de leitura eficientes.
        \item Manipular Arquivos e Diretórios: Usar os módulos \code{os} e \code{pathlib} para verificar existência, renomear e excluir arquivos.
        \item Manipular JSON e CSV: Ler e salvar dados usando os módulos nativos do Python \code{json} e \code{csv}.
    \end{itemize}

    \item \textbf{Aula 9 (\textcolor{blue}{21 de setembro de 2026} e \textcolor{blue}{25 de setembro de 2026}) - 6 horas:}

    \begin{itemize}\justifying
        \item Escopo de Variáveis: Compreender a diferença entre escopo local (interno à função) e escopo global (externo), além do uso consciente da palavra-chave \code{global}.
        \item Argumentos Padrão (Default): Criar funções com parâmetros pré-definidos (opcionais).
        \item Empacotamento de Argumentos (\code{*args}): Criar funções que aceitam múltiplos argumentos posicionais empacotados como uma tupla.
        \item Empacotamento de Argumentos Nomeados (\code{**kwargs}): Criar funções que aceitam múltiplos argumentos nomeados empacotados como um dicionário.
        \item Tipagem de Dados: Utilizar anotações de tipos (como \code{List}, \code{Dict}, \code{Tuple}) em assinaturas de funções para indicar os tipos esperados de entrada e retorno.
        \item Exercícios Práticos: Aplicar essas técnicas em exercícios práticos.
    \end{itemize}


    % ML part 2
    \item \textbf{Aula 10 (\textcolor{blue}{28 de setembro de 2026} - 3 horas:}

    \begin{itemize}\justifying
        \item Paradigma de Orientação a Objetos: Entender os conceitos de classes, objetos, atributos e métodos.
        \item O Método Construtor (\code{\_\_init\_\_}) e \code{self}: Inicializar os objetos com seus dados de forma dinâmica.
        \item Métodos de Instância: Definir comportamentos e manipular o estado interno de um objeto.
        \item Herança: Reaproveitar atributos e métodos criando classes derivadas (subclasses) a partir de uma classe base (superclass).
        \item Sobrescrita de Métodos (Overriding) e \code{super()}: Customizar funções herdadas de classes mães.
        \item Método Especial (\code{\_\_str\_\_}): Controlar como os nossos objetos são representados em formato de texto.
        \item Exercícios Práticos: Aplicar essas técnicas em exercícios práticos.
    \end{itemize}

    \item \textbf{Aula 11 (\textcolor{blue}{6 de outubro de 2026}) - 3 horas:}
    \begin{itemize}\justifying
        \item Exercícios realizados em aula.
    \end{itemize}

    \item \textbf{Aula 12:}
    \begin{itemize}
        \item Prova escrita. 
    \end{itemize}

    % matplotlib
    \item \textbf{Aula 13:}
    \begin{itemize}\justifying
        \item Introdução à Visualização de Dados: Compreender a importância da representação gráfica de dados e a arquitetura fundamental da biblioteca Matplotlib, discriminando os conceitos de \code{Figure} e \code{Axes}.
        \item Gráficos Bidimensionais (\code{plot} e \code{scatter}): Dominar a criação de gráficos de linha e de dispersão bidimensionais para a identificação de tendências e correlações em séries temporais.
        \item Customização de Elementos Gráficos: Personalizar componentes visuais essenciais, incluindo títulos, rótulos de eixos (\code{xlabel} e \code{ylabel}), legendas, grades, paletas de cores e estilos de marcadores.
        \item Histogramas e Gráficos de Barras (\code{hist} e \code{bar}): Construir representações estatísticas de distribuição de frequência e dados categóricos para análise de densidade de probabilidade de sinais.
        \item Múltiplos Subgráficos (\code{subplots}): Utilizar técnicas de particionamento de janelas gráficas para exibir múltiplos diagramas correlacionados em uma única estrutura de comparação visual.
        \item Exercícios Práticos: Aplicar essas técnicas em exercícios práticos.

    \end{itemize}


    % numpy part 1
    \item \textbf{Aula 14:}
    \begin{itemize}\justifying
        \item Entender o NumPy: Compreender por que o NumPy é tão rápido comparado às listas nativas do Python.
        \item Criar Arrays: Dominar funções de criação como \code{np.array()}, \code{np.zeros()}, \code{np.ones()}, \code{np.arange()} e \code{np.linspace()}.
        \item Inspecionar Atributos: Entender as propriedades básicas (\code{.shape}, \code{.ndim}, \code{.dtype}, \code{.size}).
        \item Fatiamento e Indexação: Acessar e recortar dados em uma e duas dimensões, incluindo o uso de máscaras booleanas.
        \item Vetorização: Realizar operações matemáticas rápidas elemento a elemento.
        \item Agregações e Eixos: Sumarizar dados utilizando \code{sum()}, \code{mean()}, \code{std()}, controlando o cálculo por linhas ou colunas (`axis`).
        \item Álgebra Linear: Dominar conceitos fundamentais como produto escalar, transposição, multiplicação de matrizes, inversão e resolução de sistemas lineares ($Ax = b$).

    \end{itemize}

    \item \textbf{Aula 15:}
    \begin{itemize}\justifying
        \item Conhecer o Pandas: Entender o propósito da biblioteca e suas duas estruturas fundamentais: \code{Series} e \code{DataFrame}.
        \item Criar Estruturas: Aprender a inicializar \code{Series} e \code{DataFrames} a partir de dicionários e listas.
        \item Importar Dados (I/O): Carregar dados externos a partir de arquivos CSV.
        \item Inspecionar Dados: Usar funções analíticas básicas como \code{.head()}, \code{.info()}, \code{.describe()} e \code{.shape}.
        \item Filtragem e Seleção: Acessar colunas, usar \code{.loc[]} / \code{.iloc[]} e aplicar filtros com máscaras booleanas.
        \item Manipulação de Dados: Criar novas colunas e tratar dados faltantes (\code{NaN}) de forma eficiente.
        \item \textbf{Agrupamento (\code{groupby})}: Sumarizar e agregar informações categóricas do dataset.
        \item Exercícios Práticos: Aplicar essas técnicas em exercícios práticos.
        \item Disponibilização da segunda lista de exercícios.
    \end{itemize}

    \item \textbf{Aula 16:}
    \begin{itemize}
        \item Entrega da primeira lista de exercícios.
        \item Resolução da primeira lista de exercícios.
    \end{itemize}

\end{itemize}



\section*{Avaliação}

A avaliação de rendimento nesta disciplina constará de três etapas distintas e independentes, compreendendo a realização de uma prova escrita individual (PI), sem qualquer tipo de consulta e realizada de forma estritamente presencial, além da entrega de duas listas (L1 e L2) de exercícios práticos, as quais também deverão ser desenvolvidas e submetidas de maneira estritamente individual. Cada um desses três instrumentos avaliativos será mensurado em uma escala de 0 (zero) a 100 (cem) pontos, de modo que a média final do discente será calculada por meio da média aritmética simples obtida a partir das três notas atribuídas ao longo do período letivo.

$$\text{MF} = \frac{(\text{PI} + \text{L1} + \text{L2})}{3}$$

As referidas listas de exercícios serão disponibilizadas de maneira gradativa ao longo do período letivo, com os respectivos prazos limites para submissão a serem oportunamente estipulados e divulgados pelo docente. Sob a perspectiva de padronização metodológica, todos os trabalhos práticos deverão ser estruturados e entregues obrigatoriamente no formato de Jupyter Notebook, devendo o discente seguir de forma rigorosa o modelo de documento e as diretrizes de formatação previamente estabelecidos e fornecidos pelo professor. O aluno estará aprovado na disciplina se $\text{MF}\geq 70,0$ pontos. Caso contrário, reprovado.


\section*{Anotações de aula}

\begin{itemize}
    \item Necessário reposição da aula do dia 1 de outubro de 2026 (2 horas).
    \item Data de prova prevista para o dia 19/23 de outubro de 2026.
    \item Atualmente com 39 horas de aula de 60 horas.
\end{itemize}


\end{document}




